\subsection{\texorpdfstring{SYSTEM OF TYPE \(G_{2}\)}{SYSTEM OF TYPE G\_2}}

\begin{enumerate}
\def\labelenumi{(\Roman{enumi})}
\tightlist
\item
  Let E be the hyperplane in \(R^{3}\) with equation
\end{enumerate}

\[
\xi_ {1} + \xi_ {2} + \xi_ {3} = 0.
\]

Let \(R\) be the set of \(\alpha \in L_0 \cap V\) such that \((\alpha | \alpha) = 2\) or \((\alpha | \alpha) = 6\) . The elements of \(R\) are

\[
\begin{array}{c} \pm (\varepsilon_ {1} - \varepsilon_ {2}), \quad \pm (\varepsilon_ {1} - \varepsilon_ {3}), \quad \pm (\varepsilon_ {2} - \varepsilon_ {3}), \quad \pm (2 \varepsilon_ {1} - \varepsilon_ {2} - \varepsilon_ {3}), \\ \pm (2 \varepsilon_ {2} - \varepsilon_ {1} - \varepsilon_ {3}), \quad \pm (2 \varepsilon_ {3} - \varepsilon_ {1} - \varepsilon_ {2}). \end{array}
\]

Then, R generates V, and \(\frac{2(\alpha|\beta)}{(\beta|\beta)}\in\mathbf{Z}\) for all \(\alpha,\beta\in R\) : this is clear if \(\beta=\pm(\varepsilon_{i}-\varepsilon_{j})\) with \(i\neq j\) ; if \(\beta=2\varepsilon_{1}-\varepsilon_{2}-\varepsilon_{3}\) for example, we have \((\alpha|\beta)\in3\mathbf{Z}\) , and again our assertion holds. Hence, R is a reduced root system in V. The number of roots is n=12.

\begin{enumerate}
\def\labelenumi{(\Roman{enumi})}
\setcounter{enumi}{1}
\tightlist
\item
  Put \(\alpha_{1} = \varepsilon_{1} - \varepsilon_{2},\alpha_{2} = -2\varepsilon_{1} + \varepsilon_{2} + \varepsilon_{3}\) . Then the roots are
\end{enumerate}

\[
\begin{array}{r l} & {\pm \alpha_ {1}, \quad \pm (\alpha_ {1} + \alpha_ {2}), \quad \pm (2 \alpha_ {1} + \alpha_ {2}), \quad \pm \alpha_ {2},} \\ & {\qquad \pm (3 \alpha_ {1} + \alpha_ {2}), \quad \pm (3 \alpha_ {1} + 2 \alpha_ {2}).} \end{array}
\]

Hence, \((\alpha_{1},\alpha_{2})\) is a basis of \(R\) . We have \(\| \alpha_{1}\|^{2} = 2,\| \alpha_{2}\|^{2} = 6.\) \((\alpha_{1}|\alpha_{2}) = -3,\) so \(R\) is a system of type \(G_{2}\) . The positive roots are \(\alpha_{1},\alpha_{1} + \alpha_{2},2\alpha_{1} + \alpha_{2},\) \(3\alpha_{1} + \alpha_{2},3\alpha_{1} + 2\alpha_{2}\) .

\begin{enumerate}
\def\labelenumi{(\Roman{enumi})}
\setcounter{enumi}{2}
\item
  We have \(h = \frac{n}{2} = 6\) .
\item
  The highest root is \(\tilde{\alpha} = 3\alpha_{1} + 2\alpha_{2} = -\varepsilon_{1} - \varepsilon_{2} + 2\varepsilon_{3}\) . We have \((\tilde{\alpha}|\alpha_{1}) = 0, (\tilde{\alpha}|\alpha_{2}) = 3\) . The completed Dynkin graph is
\end{enumerate}

\[
\alpha_ {1} \quad \alpha_ {2}
\]

\begin{enumerate}
\def\labelenumi{(\Alph{enumi})}
\setcounter{enumi}{21}
\tightlist
\item
  The inverse system is the following set of vectors:
\end{enumerate}

\[
\begin{array}{l} \pm \alpha_ {1}, \quad \pm (\alpha_ {1} + \alpha_ {2}), \quad \pm (2 \alpha_ {1} + \alpha_ {2}), \quad \pm \frac {1}{3} \alpha_ {2}, \\ \qquad \pm \frac {1}{3} (3 \alpha_ {1} + \alpha_ {2}), \quad \pm \frac {1}{3} (3 \alpha_ {1} + 2 \alpha_ {2}). \end{array}
\]

There are 10 roots not orthogonal to \(\alpha_{1}\) ; we have \(n(\beta, \alpha_{1}) = \pm 1\) for 4 of these roots, \(n(\beta, \alpha_{1}) = \pm 3\) for 4 others, and \(n(\beta, \alpha_{1}) = \pm 2\) for \(\beta = \pm \alpha_{1}\) . Hence, the squared length of \(\alpha_{1}\) for \(\Phi_{\mathbb{R}}\) is \(4(4.1 + 4.9 + 2.4)^{-1} = \frac{1}{12}\) . Hence, \(\Phi_{\mathbb{R}}(x, y) = (x|y)/24\) . We now apply formula (18) of §1, no. 12, with \(x = y = \alpha_{1}\) ; this gives

\[
2 + 4. \frac {1}{4} + 4. \frac {1}{4} = \gamma (\mathrm{R}). \frac {1}{12}
\]

so \(\gamma (\mathbb{R}) = 48\)

\begin{enumerate}
\def\labelenumi{(\Roman{enumi})}
\setcounter{enumi}{5}
\tightlist
\item
  (VI) and (VII) Half the sum of the positive roots is
\end{enumerate}

\[
\rho = 5 \alpha_ {1} + 3 \alpha_ {2}.
\]

The fundamental weights \(\omega_{1}\) and \(\omega_{2}\) are orthogonal to \(\alpha_{2}\) and \(\alpha_{1}\) , hence proportional to \(2\alpha_{1} + \alpha_{2}\) and \(3\alpha_{1} + 2\alpha_{2}\) . We have

\[
\omega_ {1} + \omega_ {2} = \rho = 5 \alpha_ {1} + 3 \alpha_ {2} = (2 \alpha_ {1} + \alpha_ {2}) + (3 \alpha_ {1} + 2 \alpha_ {2}).
\]

Hence,

\[
\omega_ {1} = 2 \alpha_ {1} + \alpha_ {2}, \quad \omega_ {2} = 3 \alpha_ {1} + 2 \alpha_ {2} = \tilde {\alpha}.
\]

\begin{enumerate}
\def\labelenumi{(\Roman{enumi})}
\setcounter{enumi}{7}
\item
  Q(R) is generated by \(\varepsilon_{1}-\varepsilon_{2}\) and \(\varepsilon_{1}-\varepsilon_{3}\) , for example. By (VI) and (VII), P(R) = Q(R). The connection index is 1.
\item
  The family of exponents has 2 terms: since 1 and \(h - 1 = 5\) are exponents, they are the only ones.
\item
  We have \((\alpha_{1},\alpha_{2}) = \frac{5\pi}{6}\) , so \(\mathrm{W}(\mathbb{R})\) is isomorphic to the dihedral group of order 12.
\item
  The only automorphism of the Dynkin diagram is the identity, so \(A(R) = W(R)\) and \(w_{0} = -1\) .
\end{enumerate}

